% !TEX root = task-2-11-6.tex
% Editable task bank. Append new BankTask blocks at the end of the document.
\documentclass[12pt,a4paper]{article}
\usepackage{physics-textbook}
\usepackage{graphicx,needspace}
\usetikzlibrary{decorations.pathmorphing}
\newif\ifShowSolutions
\ifdefined\TutorEdition
  \ShowSolutionstrue
\else
  \ShowSolutionsfalse % Change to \ShowSolutionstrue to display solutions.
\fi
\newcounter{banktask}
\newsavebox{\BankTaskBox}
\newcommand{\BankTaskContent}[3]{%
  \noindent\textbf{\thebanktask.}\quad #2\par
  \ifShowSolutions
    \Needspace{50mm}% Keep the solution heading with its opening text or diagram.
    #3\par
    \ifdefined\TutorEdition
      \ifcsname TutorSource#1\endcsname
        \nopagebreak\smallskip{\small\csname TutorSource#1\endcsname\par}%
      \fi
    \fi
  \fi
}
\newcommand{\BankTask}[3]{%
  \refstepcounter{banktask}\par\medskip
  \sbox{\BankTaskBox}{\begin{minipage}{\linewidth}%
    \BankTaskContent{#1}{#2}{#3}\end{minipage}}%
  \ifdim\dimexpr\ht\BankTaskBox+\dp\BankTaskBox\relax>\dimexpr\textheight-20mm\relax
    \Needspace{35mm}\BankTaskContent{#1}{#2}{#3}%
  \else
    \Needspace{\dimexpr\ht\BankTaskBox+\dp\BankTaskBox\relax}%
    \noindent\usebox{\BankTaskBox}\par
  \fi
  \medskip
}
\begin{document}
\BookSubsection{2.11.6}{Адиабатный процесс}


% Source 135B3C: https://ege.fipi.ru/bank/questions.php?search=1&proj=BA1F39653304A5B041B656915DC36B38&qid=135B3C
\BankTask{135B3C}{
Тонкая, упругая и мягкая оболочка герметичного воздушного шара наполнена воздухом. Как изменяются с высотой при быстром подъёме температура воздуха и концентрация молекул в шаре? Теплопроводностью оболочки шара пренебречь.

Для каждой величины определите соответствующий характер изменения:
\begin{enumerate}
\item увеличивается
\item уменьшается
\item не изменяется
\end{enumerate}

\TaskChangeTable{Температура воздуха}{Концентрация молекул}
}{
\smallskip\noindent\textbf{Решение от GPT}\par
С высотой атмосферное давление уменьшается, и воздух в шаре расширяется. При быстром подъёме \(Q=0\), поэтому \(\Delta U=-A<0\): температура уменьшается (2). Шар герметичен, число молекул \(N\) постоянно, а объём растёт, поэтому концентрация \(n=N/V\) уменьшается (2).

\(\boxed{22}\).
}

% Source 25257: https://phys-ege.sdamgia.ru/problem?id=25257
\BankTask{25257}{
В ходе адиабатного процесса над идеальным одноатомным газом была совершена работа $1495{,}8\,\text{Дж}$. В результате температура газа повысилась на $40$ °C. Найдите количество вещества этого газа.
}{
\smallskip\noindent\textbf{Решение}\par
По первому закону термодинамики $\Delta U=Q + A'.$ При адиабатном процессе Q  =  $0$. Изменение внутренней энергии для одноатомного газа $\Delta U= \frac{3}{2} v R\Delta T.$ Тогда количество вещества: 

\[v = \frac{2A'}{3R\Delta T} = \frac{2 \cdot 1495{,}8}{3 \cdot 8{,}31 \cdot 40} =3\,\text{моль}.\]

Ответ: $3$ моль.
}

% Source 50819: https://phys-ege.sdamgia.ru/problem?id=50819
\BankTask{50819}{
\begin{center}\begin{tikzpicture}[task graph,>=stealth]\draw[->](0,0)--(5,0)node[below]{\(V\)};\draw[->](0,0)--(0,3.7)node[left]{\(p\)};\node[below left]at(0,0){\(0\)};\fill(0.7,2.85)circle(1.25pt)node[above left]{\(1\)};\fill(2.75,2.85)circle(1.25pt)node[above right]{\(2\)};\fill(2.75,0.8)circle(1.25pt)node[right]{\(3\)};\draw(0.7,2.85)--(2.75,2.85);\draw[->](0.7,2.85)--(1.971,2.850);\draw(2.75,2.85)--(2.75,0.8);\draw[->](2.75,2.85)--(2.750,1.579);\draw(2.75,.8)..controls(1.7,.83)and(.77,1.4)..(.7,2.85);\draw[->](1.53,1.23)--(1.3,1.45);\end{tikzpicture}\end{center}
В качестве рабочего тела в тепловой машине используется идеальный одноатомный газ, который совершает циклический процесс, состоящий из изобарного нагревания ($1$ → $2$), изохорного охлаждения ($2$ → $3$) и адиабатного сжатия ($3$ → $1$). Отношение работы $A_{12}$, совершённой газом в изобарном процессе, к работе $A'_{31},$ совершённой над газом при адиабатном сжатии равно $2$. Определите КПД этой тепловой машины.
}{
\smallskip\noindent\textbf{Решение}\par
Проанализируем процессы. В изобарном процессе $1$–$2$ объем увеличивается. По закону Гей  — Люссака $\frac{V}{T} = const$ при увеличении объема увеличивается абсолютная температура газа. 
Работа газа в этом процессе равна $A_{12}=\nu R ( T_2 - T_1 )$ положительна, изменение внутренней энергии одноатомного газа $\Delta U_{12}= \frac{3}{2} \nu R ( T_2 - T_1 )$ положительно. Тогда по первому закону термодинамики $Q_{12}=A_{12} + \Delta U_{12}$ положительно и равно $Q_{12}= \frac{5}{2} \nu R ( T_2 - T_1 ) = \frac{5}{2} A_{12} .$ Значит, газ в этом процессе получает теплоту.
В изохорном процессе $2$–$3$ газ работу не совершает. Согласно закону Шарля $\frac{p}{T} = const$ при уменьшении давления уменьшается абсолютная температура газа. Следовательно, внутренняя энергия газа уменьшается, а по первому закону термодинамики $Q_{23}=\Delta U_{23}$ количество теплоты отрицательно, газ теплоту отдает.
В адиабатном процессе $3$–$1$ газ не получает и не отдает теплоту $Q_{31}=0,$ газ сжимается и поэтому совершает отрицательную работу, модуль которой по условию $|A_{31}|= \frac{A_{12}}{2} .$
Таким образом, работа газа за цикл равна $A=A_{12} + A_{31}= \frac{A_{12}}{2} .$ 
Количество теплоты, полученное газом за цикл $Q=Q_{12}= \frac{5}{2} A_{12} .$
Находим КПД цикла:

\[\eta= \frac{A}{Q} = \frac{2A_{12}}{2 \cdot 5A_{12}} =0{,}2=20\%.\]

Ответ: $20$\%.
}

% Source 25384: https://phys-ege.sdamgia.ru/problem?id=25384
\BankTask{25384}{
\begin{center}\begin{tikzpicture}[task graph,>=stealth]\draw[->](0,0)--(5,0)node[below]{\(V\)};\draw[->](0,0)--(0,3.7)node[left]{\(p\)};\node[below left]at(0,0){\(0\)};\fill(1,0.85)circle(1.25pt)node[below left]{\(1\)};\fill(1,3.05)circle(1.25pt)node[above left]{\(2\)};\fill(3.85,0.85)circle(1.25pt)node[below right]{\(3\)};\draw(1,0.85)--(1,3.05);\draw[->](1,0.85)--(1.000,2.214);\draw(3.85,0.85)--(1,0.85);\draw[->](3.85,0.85)--(2.083,0.850);\draw(1,3.05)..controls(1.45,1.95)and(2.55,1.58)..(3.85,.85);\draw[->](1.9,1.91)--(2.2,1.72);\end{tikzpicture}\end{center}
$1$ моль одноатомного идеального газа совершает цикл $1$–$2$–$3$–$1$, состоящий из изохоры ($1$–$2$), адиабаты ($2$–$3$) и изобары ($3$–$1$). Абсолютные температуры газа $T_{1}$  =  $400$ K, $T_{2}$  =  $600$ K, $T_{3}$  =  $510$ K. Определите КПД цикла.
}{
\smallskip\noindent\textbf{Решение}\par
КПД циклического процесса равно $\eta = \frac{A}{Q_{\,\text{пол}}} .$ Найдем работу за цикл $A=A_1{,}2 + A_2{,}3 + A_3{,}1.$ Для изохорного процесса $A_1{,}2=0.$ При адиабатном процессе из первого закона термодинамики $Q_2{,}3=A_2{,}3 + \Delta U_2{,}3$ следует, что $A_2{,}3= - \Delta U_2{,}3,$ так как $Q_2{,}3=0,$ а изменение внутренней энергии $\Delta U_2{,}3= \frac{3}{2} \nu R ( T_3 - T_2 ) .$ Для изобарного процесса $A_3{,}1=p\Delta V.$ Учитывая, что из уравнения Клапейрона  — Менделеева $pV=\nu RT,$ то $A_3{,}1=\nu R ( T_1 - T_3 ) .$ Таким образом, работа газа за цикл:

\[A=\nu R ( \frac{3}{2} ( T_2 - T_3 ) + ( T_1 - T_3 ) ) .\]

Из графика следует, что при изохорном процессе газ получал теплоту, при адиабатном процессе количество теплоты равно $0$, при изобарном процессе газ теплоту отдавал. Значит, необходимо найти количество теплоты при изохорном процессе. Из первого закона термодинамики $Q_1{,}2=A_1{,}2 + \Delta U_1{,}2,$ но при этом $A_1{,}2=0,$ а $\Delta U_1{,}2= \frac{3}{2} \nu R ( T_2 - T_1 ) ,$ то от нагревателя получена теплота $Q_{\,\text{пол}}= \frac{3}{2} \nu R ( T_2 - T_1 ) .$
Тогда КПД цикла:

\[\eta = \frac{\frac{3}{2} ( T_2 - T_3 ) + ( T_1 - T_3 )}{\frac{3}{2} ( T_2 - T_1 )} = \frac{1{,}5 \cdot ( 600 - 510 ) + ( 400 - 510 )}{1{,}5 \cdot ( 600 - 400 )} \approx 0{,}08=8\% .\]

Ответ: $8$\%.
}

% Source 9042: https://phys-ege.sdamgia.ru/problem?id=9042
\BankTask{9042}{
В вакууме закреплен горизонтальный цилиндр (см. рис.). В цилиндре находится гелий, запертый поршнем. 
\begin{center}\begin{tikzpicture}[task graph,>=stealth]\draw[fill=gray!38](0,0)rectangle(5.7,1.85);\draw[fill=white](.8,.47)rectangle(5.7,1.38);\draw[fill=orange!35](.8,.47)rectangle(3.33,1.38);\fill(3.33,.47)rectangle(3.56,1.38);\draw[fill=gray!38](3.56,.47)--(4.13,.47)--(3.56,.72)--cycle;\draw[fill=gray!38](3.56,1.38)--(4.13,1.38)--(3.56,1.12)--cycle;\draw[->](6.6,.92)--(5.72,.92);\end{tikzpicture}\end{center}
Поршень массой $90\,\text{г}$ удерживается упорами и может скользить влево вдоль стенок цилиндра без трения. В поршень попадает пуля массой $10\,\text{г}$, летящая горизонтально со скоростью $400\,\text{м/с}$, и застревает в нем. Температура гелия в момент остановки поршня в крайнем левом положении возрастает на $64$ К. Чему равно количество вещества гелия в цилиндре? Считать, что за время движения поршня газ не успевает обменяться теплом с цилиндром и поршнем.
}{
\smallskip\noindent\textbf{Решение}\par
$1$.  Запишем закон сохранения импульса:

\[m v _1= ( m + M ) v _2,\]

где m и M  — масса пули и поршня соответственно, $v _1$  — скорость пули, $v _2$  — скорость поршня с застрявшей пулей.
Поршень будет двигаться со скоростью:

\[v _2= \frac{m v _1}{m + M} .\]

$2$.  Поршень с пулей будет обладать кинетической энергией:

\[E_{\,\text{кин}}= \frac{( m + M ) v _2 ^2}{2} = \frac{m ^2 v _1 ^2}{2 ( m + M )} .\]

Изменение кинетической энергии поршня будет равно работе газа $A=\Delta E_{\,\text{к}}=E_{\,\text{к}2} - E_{\,\text{к}1}= - \frac{m ^2 v _1 ^2}{2 ( m + M )} .$ Вследствие остановки поршня $E_{\,\text{к}2}=0.$
$3$.  Поскольку газ не успевает обменяться теплом с цилиндром и поршнем, то сжатие газа будет являться адиабатичным процессом. По первому началу термодинамики $- A=\Delta U.$ Вследствие этого процесса вся механическая энергия движения поршня с пулей пойдет на нагрев газа, поэтому:

\[\Delta U= \frac{3}{2} \nu R \Delta T= \frac{m ^2 v _1 ^2}{2 ( m + M )} .\]

$4$.  Найдем отсюда количество вещества в цилиндре:

\[\nu= \frac{m ^2 v _1 ^2}{3 ( m + M ) R \Delta T} = \frac{0{,}01 ^2 \cdot 400 ^2}{3 \cdot ( 0{,}01 + 0{,}09 ) \cdot 8{,}31 \cdot 64} \approx 0{,}1\,\text{моль}.\]

Ответ: $\nu\approx 0{,}1\,\text{моль}.$
}

\end{document}
